\documentclass[10pt,a4paper]{article}
\usepackage{amsmath,amssymb,amsthm,graphicx}
\usepackage{hyperref}
\usepackage{latexsym}
\usepackage{amssymb}
\usepackage{amsthm}
\usepackage{ulem}
\usepackage{enumitem}
\setlist{itemsep=0.4pt,topsep=0pt}
\parindent 0pt
\parindent 0pt
\parskip 12pt
\newcommand{\HCF}{\mathrm{hcf}}
\newcommand{\ZZ}{\mathbb{Z}}
\newcommand{\rk}{\mathrm{rk}}
\DeclareMathOperator{\len}{length}
\DeclareMathOperator{\bl}{b}
\DeclareMathOperator{\cl}{cl}
\DeclareMathOperator{\di}{d}
\DeclareMathOperator{\Cl}{Cl}
\DeclareMathOperator{\acl}{acl}
\DeclareMathOperator{\tp}{tp}
\DeclareMathOperator{\Th}{Th}
\newcommand{\Cor}{\mathcal{C}_\mathrm{or}}
\newcommand{\Cun}{\mathcal{C}_\mathrm{un}}
\newcommand{\mons}{\mathbb{M}}
\newcommand{\strong}{\mathrm{strong}}
\newcommand{\axioms}{\Delta_{\mathrm{FSTS}}}
\newtheoremstyle{note}% hnamei
{8pt}% hSpace abovei
{14pt}% hSpace belowi
{\normalfont}% hBody fonti
{}% hIndent amounti1
{\bfseries}% hTheorem head fonti
{ }% hPunctuation after theorem headi
{.5em}% hSpace after theorem headi2
{}% hTheorem head spec (can be left empty, meaning ‘normal’)i
\theoremstyle{note}
\newtheorem{thm}{Theorem}[section]
%\newtheorem{theorem}{Theorem}
\newtheorem{lemma}[thm]{Lemma} 
\newtheorem{df}[thm]{Definition}
\newtheorem{rmk}[thm]{Remark} 
\newtheorem{prop}[thm]{Proposition}
\newtheorem{fact}[thm]{Fact}
\newtheorem{cor}[thm]{Corollary}
\newtheorem{conj}[thm]{Conjecture}
\newtheorem{ex}[thm]{Example}
\newtheorem{qu}[thm]{Question}
\newtheorem*{notation}{Notation} 
\title{A sufficiently saturated ordered system realises the criterion for every element}
\author{Silvia Barbina, Enrique Casanovas}
\date{arxiv:2411.13723v3; Creative Commons Attribution 4.0 International}
\begin{document}
\maketitle
\noindent\textit{Editorial note.} The counted statement is the author's environment at source0.tex lines 894-895 and its proof environment opens at line 896; the pairing is by adjacency and the retained proof range is the whole environment.\par\medskip
\noindent\textit{Editorial note.} The preamble is the author's own package lines, macro block and theorem-environment declarations (source lines 4-65); the author's optioned amsmath line and the fleqn-dependent layout line are not retained.\par\medskip
\noindent\textit{Editorial note.} The source carries no label-number dump and no counter values are imposed, so the printed numbering is the wrapper's default; identity is the author label, or the pinned source position where the author gave none.\par\medskip
\noindent\textit{Editorial note.} Metadata provenance: the source record of the pinned copy and the licence evidence file are both present; the licence is CC BY 4.0 with status supported, and the card's source\_version matches the value recorded in that evidence file.\par\medskip
\noindent\textit{Editorial scope.} One original result of the pinned author manuscript, counted once: in a sufficiently saturated ordered system the closure of a generated subsystem is itself (author label June\_25\_2024\_4).\par\medskip
\noindent\textit{Allowed standard prerequisites (editorial).}\begin{itemize}
\item Saturation of the ordered structure. The system is kappa-saturated as a structure in the language with the ordering, so types over small sets are realised.. It produces the element above the subsystem.
\item Binary tree above an element. In an omega-saturated ordered system there is a binary tree of height omega above any element.. It supplies the tree that the criterion of the previous proposition needs.
\item Generated subsystem. A = <A> means A is generated by itself.. It is the hypothesis on the subsystem.
\end{itemize}\medskip
\begin{cor}\label{tree4}
	Let $<$ be an HF-ordering of the infinite STS $M$ and assume %$(M,<)$ 
	$M$ is $\text{$\omega$-saturated}$ as an $(R, <)$-structure.  For every $a\in M$    there is a binary tree $T=\{a_s\mid s\in 2^{<\omega}\}$  in $M$ such that
	\begin{enumerate}
		\item $a_s>a$  for every $s$
		\item $a_s>a_t$  if  $s\subsetneq t$.
	\end{enumerate}
\end{cor}

\begin{prop} \label{May_26_2024_1} Assume  $M$ is an infinite unconfined STS, $A=\langle A\rangle \subseteq  M$  and  $a\in M\smallsetminus A$. The following are equivalent:
	\begin{enumerate}
		\item  $\cl_<(a)=\{a\}$ for some HF-ordering of $M$ with initial segment  $A$
		\item   $A\leq Aa\leq M$
		\item There are $M^\prime\succeq M$, an HF-ordering $<^\prime$ of $M^\prime$ and a binary tree $T=\{a_s\mid s\in 2^{<\omega}\}$ in $M^\prime$ such that $\cl_{<^\prime}(A)=A$,  $a=a_\emptyset$, $T>^\prime A$ and  $a_s>^\prime a_{s0},a_{s1}$ for every~$s\in$~$ 2^{<\omega}$.
	\end{enumerate}
	Moreover, these conditions imply $a\leq M$.
\end{prop}

\begin{cor}\label{June_25_2024_4}  Let $M$ be  an infinite unconfined STS  and $A\subseteq M$.  If $<$ is an HF-ordering of $M$ such that $(M,<)$ is  $\kappa$-saturated for some $\kappa>|A|$ and $\cl_<(A)=A$, then there is some $a\in M\smallsetminus A$  such that  $Aa\leq M$.
\end{cor}

\begin{proof}   By $\kappa$-saturation there is some $b\in M$  such that $b>A$. By  Corollary~\ref{tree4} we find a binary tree $T>b$ of height $\omega$ in $M$ and by the proof of $ 3 \Rightarrow 2$ in Proposition~\ref{May_26_2024_1}, the root $a$ of the tree verifies $Aa\leq M$.
\end{proof}

\end{document}
